% Pursuing Stacks (cotes 134-2 to 134-8): archive cote 134-3, site page 55 (PDF page 56), typescript page 55.
% typescript. Draft compared: ps2.tex:2778-2843.
% AI-transcribed by Claude Opus 5.5 (claude-opus-5-5, adaptive thinking, effort high), Message Batches API, batch ps1 (msgbatch_01Rv5hCaNTtUt4JHx1yqoHPL), 2026-11-12; not yet checked by eye.
\scanpage{134-3/55}
\mspage{55}
h\over{\unsure{q}}{a}s been intriguing me lately,  of understanding the ``structure'' of an
arbitrary morphism of topoi
\[ X \hand{\longrightarrow} \mathrm{Top}(B) \quad , \]
where Top(B) is the topos associated to an arbitrary (essentially small)
category B  -  a situation,  it seems,  which generalizes the situation
of a ``fibered topos'' over the indexing category B,  in the sense of
SGA 4 IV.  The construction\xout{s} of Verdier in SGA 4 V (appendix) corresponds
apparently to the case of categories B of the type $\hand{\Delta}_{/F}$  with $F \hand{\in}$ Ob$\hand{\Delta}\hat{}$,
as brought near by the standart \hand{\v{C}}ech procedure.  W\over{\unsure{n}}{h}en however the topos X
is already pretty near a topos of the type Top(B),  for instance if it is
such a topos,  to describe it's homotopy type in terms of a huge \inserted[typed]{inverse} system
of semi-simplicial objects of B$\hat{}$,  rather than just take B and keep it as
\margin{$\uparrow$}it is,  seems technically somewhat prohibitive,  at least in the present
set up (and with a distance of twenty years !).

\over{35}{36}.  Now back to test categories,  and more specifically to her majesty
$\hand{\Delta}$ -  we still have to check it is a test category indeed !  The method \inserted{suggested}
\margin{$\downarrow$}last tuesday (p. 50) does work indeed,  without any reference to ``well-
known'' facts from the semi-simplicial theory.  To check  T H 1),  we take
of course $I = \hand{\Delta}_1$ and $e_o , e_1$ as usual,  we then have for every $n \hand{\in} \underline{N}$ to con-
struct a homotopy \xout{\unsure{hx}}
\[ h_n : \hand{\Delta}_n \, x \hand{\Delta}_1 \hand{\rightarrow} \hand{\Delta}_n \quad , \]
where the product and the arrow can be interpreted as one in (Cat) or even
in (Ord) of course  (of which $\hand{\Delta}$ is a full subcategory),  and there is a
unique such h if we take for $u_a$ the ``constant'' endomorphism of $\hand{\Delta}_n$ whose
\margin{*}value is the \uline{last} element of $\hand{\Delta}_n$ \hand{(*)}. There remains only (T H 2) - namely
\margin{ok if $\exists$ final object}to prove that the objects $\hand{\Delta}_n \, x \hand{\Delta}_1$ in $\hand{\Delta}\hat{}$\xout{xxx} are aspheric $(n \hand{\geq} 0)$.  F\over{\unsure{e}}{o}r this
we cover \inserted{$\Delta'_n =$} $\hand{\Delta}_n \, x \hand{\Delta}_1$ by maximal representable subobjects,  namely maximal flags
of the ordered product set\hand{(}  these are I guess what are called the ``shuff-
les'' in semi-simplicial algebra\hand{)}. It then turns out that $\hand{\Delta'_n}$ can be obtai-
ned by successive gluing of flags,  the intersection of each flag we add
with the sup or union of the preceding ones being just a subflag,  hence
representable and aspheric.  Thus we only have to make use of the \hand{``}Mayer-
Vietoris\hand{''} type easy lemma:
\uline{Lemma}  Let  U,V two subobjects of the final object \inserted{e} of a topos \inserted{[and covering e]},  assume
U,V and U$\hand{\cap}$V aspheric,  then so is e\hand{,} i.e. so is the topos.

This finishes the proof of $\hand{\Delta}$ being a test category.  Of course,  one
cannot help thinking of the \xout{\unsure{acyclicity}} \inserted[typed]{apshericity} of \uline{all} the products $\hand{\Delta}_m \, x \hand{\Delta}_n$
as just meaning asphericity of the product ordered set (namely of the
corresponding category),  which follows for instance from the fact that it
